\begin{description}
\item[(1)] For $\Delta t_b = 2\,\mathrm{min} = 120$\,sec, the distance
  travelled by a light pulse is $R_1 = c \times \Delta t$ or
  $R_1 = 0.24$\,AU, while for $\Delta t_a = 39$\,min the maximum distance is
  $R_2 = c \times \Delta t$ or $R_2 = 4.67$\,AU. \hfill(4 points)

  Since the orbits do not intersect and the $R_2$ exceeds by far 1\,AU, the
  orbit of Seneca is exterior to that of the Earth. $R_1$ corresponds to the
  minimum relative distance of the two bodies (i.e.\ at opposition), while
  $R_2$ corresponds to the maximum relative distance (i.e.\ at conjunction).

  If $q = a\times(1-e)$ is the perihelion and $Q = a\times(1+e)$ the aphelion
  distance of Seneca, then $R_1 = q - 1$\,AU (the minimum distance of Seneca
  from the Sun minus the semi-major axis of the Earth's orbit), while
  $R_2 = Q + 1$\,AU (the maximum distance of Seneca from the Sun plus the
  semi-major axis of the Earth's orbit). \hfill(6 points)

  Thus,
  \begin{align*}
    a_{\mathrm{Sen}}(1 - e_{\mathrm{Sen}}) &= 1 + R_1 = 1.24\ \mathrm{AU} \\
    a_{\mathrm{Sen}}(1 + e_{\mathrm{Sen}}) &= R_2 - 1 = 3.67\ \mathrm{AU}
  \end{align*}
  from which one finds $a_{\mathrm{Sen}} \sim 2.46$\,AU \hfill(2 point)

  and $e_{\mathrm{Sen}} \sim 0.49$ \hfill(2 points)

  (figures rounded to 2 dec.\ digits)

\item[(2)] The period, $T_{\mathrm{Sen}}$, of Seneca's orbit can be found by
  using Kepler's 3rd law for Seneca and the Earth (ignoring their small
  masses, compared to the Sun's):
  \[ a_{\mathrm{Sen}}^{3}/T_{\mathrm{Sen}}^{2}
     = a_{\mathrm{Earth}}^{3}/T_{\mathrm{Earth}}^{2} = 1 \]
  \hfill(6 points)

  from which one finds $T_{\mathrm{Sen}} \sim 3.87$\,yr \hfill(2 point)

  (Alternatively one can use Kepler's law for Seneca only and use natural
  units for the mass of the Sun and $G$, i.e.\ $[\mathrm{mass}] =
  1\,M_{\mathrm{Sun}}$, $[t] = 1$\,yr and $[r] = 1$\,AU, in which case
  $T_{\mathrm{Sen}} = a_{\mathrm{Sen}}^{3/2} = 3.86$\,yr) \hfill(8 points)

  Assuming non-retrogate orbit, % sic: "non-retrogate" in source
  the synodic period, $T_{\mathrm{syn}}$ of Seneca and Earth is given by
  \[ 1/T_{\mathrm{syn}} = 1/T_{\mathrm{Earth}} - 1/T_{\mathrm{Sen}}
     \quad (\text{Seneca is superior}) \]
  \hfill(6 points)

  which gives $T_{\mathrm{syn}} = 1.35$\,yr \hfill(2 point)

\item[(3)] From Kepler's law, we can calculate the mass of the Sun
  \[ 4\pi^2\, a_{\mathrm{Sen}}^{3}/T_{\mathrm{Sen}}^{2} = G\,M_{\mathrm{Sun}},
     \quad\text{i.e.}\quad
     M_{\mathrm{Sun}} \sim 1.988\times10^{30}\ \mathrm{kg}
     \quad(\text{depending on accuracy}) \]
  \hfill(4 points)

  Then, using Kepler's 3rd law for Seneca and Jupiter, one gets
  \[ \left(a_{\mathrm{Jup}}^{3}/T_{\mathrm{Jup}}^{2}\right) \big/
     \left(a_{\mathrm{Sen}}^{3}/T_{\mathrm{Sen}}^{2}\right)
     = \left(M_{\mathrm{Sun}} + M_{\mathrm{Jup}}\right)/M_{\mathrm{Sun}}
     = 1 + x \]
  \hfill(4 points)

  where $x$ is the mass ratio of Jupiter to the Sun. Then, solving for $x$ we
  get (depending on the accuracy used in determining the elements of Seneca's
  orbit)
  \[ x = 0.0016, \quad\text{thus}\quad
     M_{\mathrm{Jup}} = 3.2\times10^{27}\ \mathrm{kg} \]
  \hfill(4 points)
\end{description}
